\section{The limiting topos and the information retained by a germ}

Throughout this part, a covering of a topological space means a map
locally isomorphic over an open neighbourhood $U$ to $U\times A\to U$
for a discrete set $A$, with no finiteness restriction unless one is
stated. Its corresponding sheaf is a locally constant object. A limit
of topoi means a bicategorical
limit for geometric morphisms. An intersection of sublocales is not,
without proof, identified with an intersection of sets of points.
In a classifying topos of a nondiscrete topological group, a continuous
action need not itself be a locally constant object; F.10 makes this
distinction explicit.

\begin{statement}{F.1}{Open}{The arbitrary covering-germ comparison}
The literal inverse-topos limit is determined for every filter of open
subspaces by F.7 below. We also determine the comparison for a point of a
manifold, for a closed submanifold with a cofinal tubular system, and for
its residual germ. We do not determine, for an arbitrary space and filter,
when every possibly infinite covering and every morphism of coverings on
the intersection sublocale descends to one common neighbourhood. In
particular, local connectedness and semilocal simple connectedness of the
individual stages do not justify the proposed pro-groupoid replacement
in Lead 14: F.8 gives a counterexample.
F.11 gives a compact closed subset for which the ordinary comparison
already fails on morphisms of infinite coverings. F.14 gives an exact
stability criterion when the limit and stages are connected covering
spaces and the induced fundamental-group maps are surjective.
F.17 classifies all Hausdorff point germs, and F.18 classifies the
residual germs of all closed subsets of regular Hausdorff spaces.
\end{statement}

\begin{statement}{F.2}{Known}{Finite covering descent in a coherent limit}
Let $(\mathcal E_i)$ be a small cofiltered diagram of coherent
Grothendieck topoi and coherent geometric morphisms, and let
$\mathcal E=\operatorname*{2\!\varprojlim}_i\mathcal E_i$.
Restriction induces an equivalence
\[
 \operatorname*{2\!\varinjlim}_i\operatorname{FLC}(\mathcal E_i)
       \simeq \operatorname{FLC}(\mathcal E),
\]
where $\operatorname{FLC}$ denotes finite locally constant objects.
If $X$ is smooth over a characteristic-zero field, $\bar x$ is a geometric
point on a strict normal crossings divisor with $r\geq1$ local branches,
and $R=\mathcal O^{\mathrm{sh}}_{X,\bar x}$, then the limit of the small
\'etale topoi of the complements in pointed affine \'etale
neighbourhoods is
\[
 \bigl(\operatorname{Spec}R[1/(t_1\cdots t_r)]\bigr)_{\mathrm{et}}.
\]
Its finite covering group, relative to a compatible geometric basepoint,
is $\widehat{\mathbf Z}(1)^r$. For a closed geometric point on a smooth
curve this entire topos is $B\widehat{\mathbf Z}(1)$. For a regular local
ring of dimension greater than one it is not the classifying topos of
that group.
\end{statement}
\begin{proof}
SGA 4, Expos\'e VI, Theorems 8.2.3 and 8.3.13 identify the limit topos
with the sheaves on the filtered colimit of the coherent sites and
identify its coherent objects with the filtered colimit of the coherent
objects of the stages \cite{fghSGA4}. Here is the extra step from coherent
to finite locally constant. A finite locally constant object is coherent:
a trivializing cover may be refined by coherent objects, then reduced to
a finite cover because the terminal object is coherent; over each member
it is a finite coproduct of that coherent object. The object, the finite
trivializing cover, the finitely many trivializations, their inverses,
and their equations therefore descend to a common stage. The covering
property descends too: the finite coherent covering families generate
the colimit-site topology, as proved in VI 8.3.6. Thus the descended
object becomes finite locally constant at a later stage. Full faithfulness
is inherited from the equivalence on coherent objects. This argument
uses finite trivializations; it makes no assertion about infinitely many
simultaneous trivializations.

For the scheme assertion, affine \'etale algebras of finite presentation,
their maps, and finite covering families descend through an affine
inverse limit. In the present diagram the limit ring is $R$, and
localization commutes with its filtered colimit. This proves the stated
identification of sites and hence topoi; the exact site statement is
Stacks Project, Lemma 59.51.2, tag 0EYL, and finite locally constant
descent is also Lemma 59.73.11, tag 0GL2 \cite{fghStacksLimits}.

We spell out the ramification argument used to identify the covering
group. A connected finite \'etale cover of the complement has finite
normalization over $R$, since $R$ is excellent. Its only possible
ramification divisors are $t_i=0$, and all their ramification indices
are prime to the residue characteristic, which is zero. Choose a common
multiple $n$ of these indices. Adjoining $n$th roots of all the $t_i$
kills this ramification by Abhyankar's lemma. The resulting normalization
is \'etale everywhere by purity of the branch locus. The root-adjoined
local regular ring is strictly henselian, so the normalization splits.
Conversely the covers obtained by adjoining roots of the $t_i$ have the
independent deck transformations $\mu_n^r$. Thus their inverse limit is
$\widehat{\mathbf Z}(1)^r$. This use of Abhyankar and purity is the usual
strict normal crossings form of SGA 1, Expos\'e XIII, Corollary 5.3
\cite{fghSGA1}.

If $\dim R=1$, the localized ring is its fraction field. Every finite
separable extension is a Kummer extension of its uniformizer: after
removing the tame ramification index the residual extension is unramified
and therefore trivial. The small \'etale topos of a field is the category
of continuous sets for its absolute Galois group. If $\dim R>1$, the
localized ring has a proper nonzero prime. For $r=1$ use the prime
$(t_2)$, extending $t_1$ to a regular system of parameters; for $r>1$ use
$(t_1+t_2)$. These primes avoid $t_1\cdots t_r$. Its spectrum therefore
has a proper nonempty open subset, giving a third subobject of the
terminal sheaf. A classifying topos of a group has only the empty and
terminal subobjects of its terminal object. This proves the last
assertion without confusing finite coverings with all sheaves.
\end{proof}

\begin{statement}{F.3}{Open}{The two readings of the manuscript parenthesis}
For the open-subspace reading of \S53.5, cote 140-4, folios 768'--769,
scans 242--243, F.7 gives the literal limit and F.9 gives the exact
replacement for a normal crossings puncture. For the coherent \'etale
reading F.2 proves finite covering descent. Leroy's thesis concerns the
fundamental topos generated by locally constant objects and a
van Kampen theorem; its exact assertions are recorded in F.13. These
results do not settle the requested unrestricted classification of all
filters and all infinite covering germs in F.1.
\end{statement}

\begin{statement}{F.4}{Proved}{Tangent tori and cuspidal homotopy fixed points}
Let $x$ be a point on a strict complex normal crossings divisor $D$ in a
complex manifold $X$, with $r$ branches at $x$. For a decreasing cofinal
sequence of open neighbourhoods $V_n$ of $x$,
\[
 \operatorname*{holim}_n\operatorname{Sing}(V_n\setminus D)
       \simeq \operatorname{Sing}((S^1)^r).
\]
The right side is the fibre over $x$ of the real oriented blow-up. This
is a weak equivalence of homotopy types; the ordinary intersection of
the punctured neighbourhoods is empty.

Let $U/L$ be a smooth geometrically connected hyperbolic curve over a
number field, with an $L$-rational cusp $c$. The geometric punctured
strictly henselian neighbourhood has profinite type $B I_c$, where
$I_c=\widehat{\mathbf Z}(1)$. After choosing one tangential splitting,
its continuous homotopy fixed-point space has
\[
 \pi_0=H^1_{\mathrm{cts}}(G_L,I_c)
       =\varprojlim_n L^\times/(L^\times)^n,
 \qquad \pi_j=0\quad(j>0).
\]
Without the splitting, the first displayed set is naturally a torsor
under that group. Its map to the section classes of $\pi_1(U/L)$ is
injective, with image the packet at $c$. Rational nonzero tangent
vectors give the dense, proper $L^\times$-subtorsor.
\end{statement}
\begin{proof}
Replace the cofinal sequence by coordinate polydiscs centred at $x$,
using the cofinality invariance of homotopy inverse limits. The
complements have coordinates
$(\Delta^*)^r\times\Delta^{d-r}$. The angular projection to $(S^1)^r$
is compatible with all radial inclusions and is a homotopy equivalence
at every stage: its radial and unpunctured-coordinate fibres are
contractible. Applying singular complexes and the homotopy-limit
functor gives the equivalence to the homotopy limit of the constant
torus diagram. A compatible tangent ray determines a point of this
homotopy limit, with radial paths providing the coherence; it need not
be a point of the empty strict inverse limit.

In the arithmetic case, a continuous homotopy section of
$B(I_c\rtimes G_L)\to BG_L$ is a continuous splitting modulo
$I_c$-conjugation. In a chosen splitting it is precisely a continuous
one-cocycle. An automorphism of such a section is an invariant element
of $I_c$, since $I_c$ is abelian; there are no higher homotopies for
these one-types. For every prime $\ell$, the cyclotomic character of a
number field has infinite image in $\mathbf Z_\ell^\times$, so
$\mathbf Z_\ell(1)^{G_L}=0$. Consequently $I_c^{G_L}=0$.

Choose the cofinal sequence of coefficient moduli $n!$. The inverse
system of groups $H^0(G_L,\mu_{n!})$ consists of finite groups and
is Mittag--Leffler. Its first derived inverse limit vanishes. The
continuous-cohomology inverse-limit exact sequence therefore gives
\[
 H^1_{\mathrm{cts}}(G_L,I_c)
       =\varprojlim_n H^1(G_L,\mu_{n!}).
\]
Kummer theory identifies the terms with
$L^\times/(L^\times)^{n!}$. This is the precise reason that no
$\varprojlim^1$ term occurs here; it is not a general assertion about
integral cohomology in degree two.

Stix, \emph{On cuspidal sections}, Theorem 14, proves packet
injectivity and disjointness from rational-point sections and other
cusps, assuming point-section injectivity for proper genus-at-least-two
curves over finite extensions \cite{fghStix}. For number fields that
hypothesis follows from the Jacobian Kummer map and the Mordell--Weil
theorem: the difference of points with equal section classes is divisible
by every integer in the finitely generated group of rational Jacobian
points, and is therefore zero. Passing to finite \'etale covers of the
section gives the same argument for distinct lifts of a cusp. This
checks exactly the hypothesis needed in Stix's theorem.

Scaling a rational tangent vector by $a\in L^\times$ changes the
splitting by the Kummer class of $a$, as one sees on the parameter
covers $t\mapsto t^{1/n}$. The map from $L^\times$ to its completion
is injective: all valuations first vanish for an infinitely divisible
element, and the finitely generated unit group then gives the claim.
Its image is dense because each finite set of modulus conditions is
represented at their common multiple. It is proper because $L^\times$
is countable whereas the closure of the powers of a single element
with a nonzero finite-place valuation contains a copy of
$\widehat{\mathbf Z}$ and is uncountable.
\end{proof}

\begin{statement}{F.5}{Open}{The proposed homotopy reformulation of the section conjecture}
For a hyperbolic curve in characteristic zero, the condition that every
section is a homotopy fixed point of its geometric \'etale homotopy
type is automatic. Hence adding the alternative that it be a homotopy
fixed point of a punctured cusp neighbourhood does not establish the
section conjecture. The desired nonautomatic assertion is surjectivity
of
\[
 U(L)\;\amalg\!\coprod_{c\in(\overline U\setminus U)(L)}
       \operatorname{Sec}(D_c/G_L)
       \longrightarrow
       \pi_0\bigl(\operatorname{Et}(U_{\overline L})^{hG_L}\bigr),
\]
where each section set on the left is modulo its geometric inertia.
This is the section conjecture with cuspidal packets. Its general
surjectivity over number fields is not proved here. In particular the
requested identification of the entire first homotopy fixed-point set
with $U(L)$ is false whenever $U$ has a rational cusp.
\end{statement}
\begin{proof}[Proof of the correction]
The geometric \'etale homotopy type of a hyperbolic curve in
characteristic zero is a profinite $K(\pi,1)$; the curve statement is, for example,
Schmidt--Stix, Lemma 2.7(a) \cite{fghSchmidtStix}. Thus the classifying-space
construction identifies the homotopy fixed-point groupoid with the
groupoid of splittings of the arithmetic fundamental-group extension.
This is the same elementary description of maps over $BG_L$ used in
F.4, with a nonabelian geometric group. Every section therefore lies
in the first alternative of the proposed condition, without a
geometric-origin assertion. By F.4, any rational cusp gives a section,
and Stix's disjointness theorem says that this section is not the
section of an element of $U(L)$. The two proposed sets therefore cannot
be identified by the natural point-section map. The displayed corrected
map is injective by the same point and packet injectivity and is
surjective exactly when every section has one of those geometric
origins. This proves the reformulation while leaving its general
surjectivity undecided.
\end{proof}

\begin{statement}{F.6}{Refuted}{The unrestricted stack version of Lead 21}
Lead 21 is false when its torus formula is applied without retaining
stack isotropy. For the smooth \'etale analytic stack
$[\Delta/C_2]$ with trivial action and divisor $[\{0\}/C_2]$, the
punctured homotopy germ and tangent fibre have homotopy type
$BC_2\times S^1$, whose first integral homology is
$C_2\oplus\mathbf Z$, not $\mathbf Z$.
Here the tangent fibre means the inverse image of the residual gerbe
of the underlying stack point. The geometric 2-fibre over a chosen
map $*\to[\Delta/C_2]$ is $S^1$. The inverse-topos limit of intrinsic
closed blow-up neighbourhoods instead retains $BC_2\times S^1$;
so (B) is also false if its right side is read as that geometric
2-fibre. It is correct with the residual gerbe, under F.16's
closed-stage hypotheses.
For manifolds, the correct assertions (A), (B), and (C), with their
hypotheses made explicit, are proved in F.15, F.9, and F.12. For an
\'etale stack the corresponding assertions use its actual tangent
stack and require compatible intrinsic neighbourhoods; chartwise
shrinking alone does not supply that compatibility.
\end{statement}
\begin{proof}
The classifying space of the product quotient is the product of the
classifying space of $C_2$ with the disc. Removing the divisor and
shrinking the disc changes the second factor to $S^1$ up to homotopy
and leaves the first factor unchanged. Abelianization of
$C_2\times\mathbf Z$ gives the stated first homology group. The
finite torsion factor is checked by the bar-complex computation in
\texttt{code/fgh\_topology\_checks.py}; the proof does not identify
the coarse-space tangent circle with the tangent stack.

The same counterexample has an effective orbifold realization.
Take $[\Delta_z\times\Delta_w/C_2]$, where the nontrivial element
acts by $(z,w)\mapsto(-z,w)$, and take the divisor $w=0$ at
the origin. The punctured bidisc retracts equivariantly onto
$\{z=0\}\times S^1$. The action there is trivial, so its quotient
homotopy type is again $BC_2\times S^1$. Thus excluding ineffective
stabilizers would not repair the asserted torus formula.
\end{proof}

\begin{statement}{F.7}{Proved}{Localic intersection, and closed germs in a regular space}
Let $S$ be a space and $(S_i)$ a filter basis of open subspaces. In the
locale of $S$, let $L=\bigcap_i S_i$ be the meet of the associated open
sublocales. Then
\[
 \operatorname*{2\!\varprojlim}_i\operatorname{Sh}(S_i)
       \simeq\operatorname{Sh}(L).
\]
If $S$ is regular Hausdorff, $Y\subset S$ is closed, and the $S_i$
are cofinal among neighbourhoods of $Y$, then $L=Y$ as a closed
sublocale. Moreover,
\[
 \operatorname*{2\!\varprojlim}_i\operatorname{Sh}(S_i\setminus Y)
       \simeq\operatorname{Sh}(\varnothing).
\]
In particular the ordinary germ of a manifold at a point has limit
$\operatorname{Set}$, while its nonempty residual germ has degenerate
limit.
\end{statement}
\begin{proof}
An open subtopos is a sublocale, and a geometric morphism into
$\operatorname{Sh}(S)$ factors through $\operatorname{Sh}(S_i)$
exactly when its inverse-image functor takes the subterminal $S_i$
to the terminal object. Requiring this for every $i$ is the universal
property of the intersection sublocale. It gives the same universal
property for a compatible cone of geometric morphisms, including
2-morphisms, and proves the first assertion.

For the second assertion, let $g:\mathcal T\to\operatorname{Sh}(S)$
be such a cone. For each point $z\notin Y$, regularity provides an open
$W_z$ with $z\in W_z$ and $\overline{W_z}\cap Y=\varnothing$.
Some $S_i$ lies in $S\setminus\overline{W_z}$. Hence
$g^*W_z=g^*(W_z\cap S_i)=0$. These opens cover $S\setminus Y$,
so $g^*(S\setminus Y)=0$. This is exactly the factorization criterion
through the closed subtopos $\operatorname{Sh}(Y)$. Conversely its
inclusion factors through every $S_i$. This proves the equality of
sublocales, rather than just their point sets. In the residual diagram
the same $W_z$ form an open cover of a fixed ambient space
$S\setminus Y$, and every member pulls back to zero in every cone.
Thus the terminal object of $\mathcal T$ is zero. Such a cone exists
only from a degenerate topos, proving the last assertion.
\end{proof}

\begin{statement}{F.8}{Refuted}{Some covering germs are not the coverings of any Grothendieck topos}
The final assertion of Lead 14 fails even when every stage is
discrete. Put $S=\mathbf N$ with the discrete
topology and $S_n=\{n,n+1,\ldots\}$. Then every $S_n$ is locally
connected and semilocally simply connected. Nevertheless
\[
 \operatorname*{2\!\varprojlim}_n B\Pi_1(S_n)
       =\operatorname{Sh}(\varnothing),
 \qquad
 \operatorname*{2\!\varinjlim}_n\operatorname{Cov}(S_n)
       \text{ is nondegenerate}.
\]
The latter category is not the category of locally constant objects
of \emph{any} Grothendieck topos: it has no countable coproduct of
its terminal object, whereas every such locally constant category
contains the constant object $\coprod_{n\geq1}1$.
The germ category is also not itself a Grothendieck topos. Its subobjects
of the terminal object form $\mathcal P(\mathbf N)/\mathrm{Fin}$,
which is not a complete lattice. Thus the stagewise local hypotheses
in Lead 14 do not imply even the claimed existence of a replacement
Grothendieck topos with precisely these coverings.
\end{statement}
\begin{proof}
Each fundamental groupoid here is the discrete groupoid on $S_n$;
its classifying topos is $\operatorname{Set}^{S_n}=
\operatorname{Sh}(S_n)$. Every singleton is absent at a later stage,
so any cone pulls each singleton subterminal to zero. Since the
singletons cover $S_1$, the cone's terminal object is zero. The inverse
limit is degenerate. In the filtered covering category the empty and
one-sheeted covers never become isomorphic, since all $S_n$ are
nonempty. We prove the stronger failure of a countable coproduct.
Suppose that an object $A=(A_n)$ and maps $i_j:1\to A$ had that
universal property. Discarding a finite initial segment, all $A_n$
are nonempty, and each $i_j$ can be represented by a sequence
$a_j(n)\in A_n$, filling its finitely many undefined entries
arbitrarily. For any two different $j,k$, the coproduct property
provides a morphism $A\to2$ taking $i_j$ to $0$ and $i_k$ to $1$.
Therefore $a_j(n)\ne a_k(n)$ for all sufficiently large $n$.
Consequently every finite list $a_1(n),\ldots,a_m(n)$ is pairwise
distinct for all sufficiently large $n$.

Choose integers $k(n)\to\infty$ such that
$a_1(n),\ldots,a_{k(n)+1}(n)$ are distinct. Set
$b_n=a_{k(n)+1}(n)$, and let $h_n:A_n\to\{0,1\}$ be the
characteristic function of $\{b_n\}$. These define a nonzero
germ morphism $h:A\to2$, since each $h_n$ takes the value $1$.
For every fixed $j$, however, $h_n(a_j(n))=0$ eventually, since
$k(n)\geq j$ eventually. Thus $hi_j=0$ for every $j$, in
contradiction to uniqueness in the coproduct property. In a
Grothendieck topos the countable coproduct of its terminal object
is a constant, hence locally constant, object. Its universal
property holds in the full subcategory of locally constant
objects as well. An equivalence with that subcategory is therefore
impossible.

A subobject of the germ category's terminal object is a subset of a tail,
and two such subsets determine the same subobject exactly when they
differ by finitely many points. This gives the stated Boolean algebra.

To prove noncompleteness, partition $\mathbf N$ into countably many
infinite subsets $A_j$. If $[B]$ were the join of the $[A_j]$, then
$A_j\setminus B$ would be finite for each $j$. Choose
$b_j\in B\cap A_j$ for each $j$. The set
$B'=B\setminus\{b_j:j\geq1\}$ is another upper bound, but
$[B']<[B]$. Thus there is no join. Every Grothendieck topos has a
complete subobject lattice, proving the additional claim. Neither
argument treats a filtered category of finite coverings as if it
were already its Ind-completion.
\end{proof}

\begin{statement}{F.9}{Proved}{The normal crossings germ and compact blow-up stages}
For the punctured coordinate neighbourhoods of a strict complex normal
crossings divisor at a point with $r$ branches, the covering-germ
category is $\mathbf Z^r$-sets. A topos encoding precisely these
coverings is $B\mathbf Z^r$. Its finite covering category is equivalent
to that of $B\widehat{\mathbf Z}(1)^r$ after choosing an orientation.
The literal sheaf-topos limit is degenerate, as in F.7.
If $N_n$ are compact neighbourhoods of $x$ in one fixed compact
coordinate chart, nested with intersection $\{x\}$, then in the real
oriented blow-up $\tau$,
\[
 \operatorname*{2\!\varprojlim}_n\operatorname{Sh}(\tau^{-1}N_n)
       \simeq\operatorname{Sh}(\tau^{-1}(x)).
\]
For a closed smooth submanifold with a cofinal tubular system the
ordinary covering germ is instead its own covering category.
\end{statement}
\begin{proof}
The angular deformation retractions in F.4 identify every restriction
functor on coverings with the identity on $\mathbf Z^r$-sets. For
finite sets an action of $\mathbf Z^r$ extends uniquely to a continuous
action of its profinite completion. Tangential fibres are evaluation
along rays; changing a connecting angular path changes transport by
its integral winding vector. On finite fibres the corresponding
automorphism torsors complete from $\mathbf Z^r$ to
$\widehat{\mathbf Z}(1)^r$.

For the compact assertion the spaces $\tau^{-1}N_n$ are closed in
one compact Hausdorff space. A cone into their sheaf topoi pulls the
open complement of every closed stage to zero. It therefore pulls
their union to zero, and factors through the closed subtopos of their
intersection. Conversely that intersection factors through every
stage. This proves the claimed bicategorical universal property.
Compactness ensures here that all the chosen stages are closed in the
fixed Hausdorff ambient space; the closed-subtopos argument itself
does not require a countable indexing set. The ordinary
tubular-neighbourhood assertion follows by the
compatible deformation retractions to the submanifold, which induce
equivalences on fundamental groupoids and hence on arbitrary covering
categories. Neither argument extends chartwise over a nonproper
stack presentation without an intrinsic common stage.
\end{proof}

\begin{statement}{F.10}{Proved}{A precise prodiscrete replacement for surjective connected towers}
Let $G_1\leftarrow G_2\leftarrow\cdots$ be a countable tower of
surjective homomorphisms of discrete groups. Give
$G=\varprojlim_nG_n$ its inverse-limit topology, and write $BG$ for
the topos of discrete continuous $G$-sets. Then
\[
 BG\simeq\operatorname*{2\!\varprojlim}_n BG_n,
 \qquad
 \operatorname{LC}(BG)\simeq
       \operatorname*{2\!\varinjlim}_n(G_n\text{-}\operatorname{Set}).
\]
More explicitly, a continuous $G$-set is locally constant as an
object of $BG$ if and only if one kernel
$\ker(G\to G_n)$ acts trivially on the whole set. Thus, if a
pointed connected germ has stagewise covering classifications by this
surjective tower, $BG$ realizes its covering-germ category. This
assertion does not require the discrete groups $G_n$ to be finite.
\end{statement}
\begin{proof}
The projections $G\to G_n$ are surjective: starting from an element
at one stage, recursively choose lifts at every later stage.
Their kernels form a neighbourhood basis of the identity. Inflation
therefore identifies $G_n$-sets with a full subcategory of $BG$ and
preserves finite limits and all colimits, so it is the inverse image
of a geometric morphism $BG\to BG_n$.

Every transitive continuous $G$-set is $G/U$ for an open subgroup
$U$. Some $\ker(G\to G_n)$ is contained in $U$, so that transitive
set comes from $BG_n$. The union of these full subcategories thus
generates $BG$. Leroy's Proposition 3.3.3, whose hypotheses are
exactly these full geometric inclusions and the generating union,
identifies $BG$ with the indicated 2-limit \cite{fghLeroy}.

If a kernel acts trivially on $S$, the epimorphism
$G/\ker(G\to G_n)\to1$ trivializes $S$; hence $S$ is locally
constant. Conversely choose an epimorphic family trivializing a
locally constant $S$, a nonempty member $T$, and an element $t\in T$.
Its stabilizer $U$ is open. Restrict the trivialization over $T$ to
the fibre above $t$: it says that the $U$-action on the whole set $S$
is trivial. Some kernel is contained in $U$, proving the criterion.
Full faithfulness follows from surjectivity of the projections, and
the criterion proves essential surjectivity of the displayed functor.

For example $\coprod_{n\geq1}\mathbf Z/2^n$ is a continuous
$\mathbf Z_2$-set. It is not locally constant in $B\mathbf Z_2$:
no open subgroup acts trivially on every summand. Thus replacing
``locally constant objects'' by ``all objects'' in the second
equivalence would be false. This is precisely the distinction between
Leroy's $\operatorname{LC}$ and its sum-completion
$\operatorname{SLC}$.
\end{proof}

\begin{statement}{F.11}{Refuted}{Compactness alone does not give descent for arbitrary coverings}
Let $C\subset[0,1]\subset\mathbf R$ be the middle-thirds Cantor
set, and let $\Phi$ be its filter of open neighbourhoods in
$\mathbf R$. The literal topos limit is $\operatorname{Sh}(C)$,
but the comparison
\[
 \operatorname*{2\!\varinjlim}_{V\in\Phi}\operatorname{Cov}(V)
          \longrightarrow\operatorname{Cov}(C)
\]
is not full. Its failure occurs between two trivial coverings, one
with countably infinite fibre and the other with two-point fibre.
For these same two coverings every morphism using only finitely many
of the specified sheets does extend to a common neighbourhood.
\end{statement}
\begin{proof}
Identify $C$ with $\{0,1\}^{\mathbf N}$ by its ternary expansion
using only digits $0$ and $2$, and let $d_n:C\to\{0,1\}$ be
the $n$th coordinate. The formula
\[
  (c,n)\longmapsto(c,d_n(c))
\]
defines a continuous morphism of coverings
$C\times\mathbf N\to C\times\{0,1\}$: each sheet is open,
and every $d_n$ is locally constant.

If this morphism extended between the corresponding trivial covers
on an open neighbourhood $V$ of $C$, each of its sheet functions
would be constant on each component interval $J$ of $V$.
Consequently every $d_n$ would be constant on $J\cap C$.
The coordinates separate points of $C$, so $J\cap C$ could contain
at most one point. But $C$ is perfect, and any open interval
containing a point of $C$ contains a second point of $C$. This is a
contradiction. A morphism in the filtered colimit between the two
fixed trivial objects is represented by a morphism on some such $V$,
so the contradiction is exactly failure of fullness.

Any finite list of digits is constant on each of finitely many
clopen cylinder subsets of $C$. These compact subsets have disjoint
open neighbourhoods in $\mathbf R$, on which the finitely many
functions extend simultaneously. The obstruction is therefore the
absence of one common neighbourhood for the infinite family, rather
than failure of the separate finite extension problems. F.7 computes
the literal limit because $C$ is closed in a regular Hausdorff space.
\end{proof}

\begin{statement}{F.12}{Proved}{The singular-chain derived limit records the missing torus}
Let $C_n=C_\bullet(V_n\setminus D;\mathbf Z)$ for a cofinal sequence
of coordinate polydiscs as in F.4. Then, degree by degree,
$\varprojlim_n C_n=0$. If $Q_\bullet$ is the degreewise first derived
inverse limit, then
\[
 R\!\varprojlim_n C_n\simeq Q_\bullet[-1],
 \qquad H_{j+1}(Q_\bullet)\simeq
       H_j((S^1)^r;\mathbf Z).
\]
Here $[-1]$ is the homological shift with
$H_j(Q[-1])=H_{j+1}(Q)$. The analogous assertion for a stack uses
its actual homotopy fibre, including isotropy.
\end{statement}
\begin{proof}
Regard each singular chain group as a subgroup of the chain group of
the first punctured neighbourhood. A chain has finite support, and
each simplex has compact image disjoint from $D$ at the base point.
In particular no simplex can have image contained in every shrinking
punctured neighbourhood: their intersection is empty. Linear
independence of the singular-simplex basis gives the zero inverse
limit in each degree.

For a countable tower of abelian groups the two-term Roos complex
$\prod C_n\xrightarrow{1-\mathrm{shift}}\prod C_n$ computes the
derived inverse limit. Its degreewise kernel is zero and its cokernel
is $Q_\bullet$, proving the first quasi-isomorphism. The Milnor exact
sequence for its homology is
\[
 0\longrightarrow\varprojlim{}^1 H_{j+1}(C_n)
 \longrightarrow H_j(R\!\varprojlim C_n)
 \longrightarrow\varprojlim H_j(C_n)\longrightarrow0.
\]
The homology systems are constant up to the compatible angular
isomorphisms, so the first term vanishes and the last is the torus
homology. Combining the two formulas gives the claimed degree shift.
No interchange of chain inverse limits with homology was used.
\end{proof}

\begin{statement}{F.13}{Known}{What Leroy proves, and what it does not identify}
For a locally connected Grothendieck topos $\mathcal E$, Leroy's
Theorem 2.4 constructs the topos
$\operatorname{SLC}(\mathcal E)$ of sums of locally constant
objects and the geometric morphism
$\mathcal E\to\operatorname{SLC}(\mathcal E)$ whose inverse
image is the inclusion. Its Galois objects generate it. Definition
3.4.1 and the subsequent remark identify this as the fundamental
locally Galois topos. Propositions 3.3.2--3.3.4 present it by
groupoid topoi with fully faithful inverse-image transitions and
specified trivializing covering sieves. Proposition 4.1.1 constructs
a 2-\emph{colimit} of topoi by cartesian sections; the van Kampen theorem in \S4.5
compares the associated fundamental topoi when the fibre topoi are
locally connected.

These are statements about a chosen topos and its locally constant
objects, or about that stated colimit. They do not identify the
literal inverse limit of all sheaf topoi of every punctured filter
with the filter's covering category. F.8 and F.11 exhibit two distinct
obstructions; F.9 and F.10 give precise replacement results.
\end{statement}
\begin{proof}[Source verification and application]
The original scan was opened at printed pages 15--16 for Theorem
2.4 and pages 26--29 for Propositions 3.3.1--3.3.3. The provisional
transcription was read in chapters 2--4 and checked against those
scan pages. The source and transcription locations are recorded in
the bibliography \cite{fghLeroy}. In particular the fully faithful
transition and generating-union hypotheses of 3.3.3 are explicit;
F.10 verifies them, whereas the disappearing discrete components of
F.8 do not. For a normal crossings puncture, the stable fundamental
groupoid already gives $B\mathbf Z^r$, as proved directly in F.9.
This resolves the parenthesis in these specified readings without
attributing an unrestricted continuity theorem to Leroy.
\end{proof}

\begin{statement}{F.14}{Proved}{An exact stability criterion under surjective fundamental-group maps}
Suppose the limit sublocale is a connected, locally path connected,
semilocally simply connected space $Y$, and the filter has a cofinal
system of connected, locally path connected, semilocally simply
connected stages $S_i$ containing $Y$. Choose a point $y\in Y$.
Assume every map
\[
 \phi_i:\pi_1(Y,y)\longrightarrow\pi_1(S_i,y)
\]
is surjective. Then the covering comparison
$\operatorname*{2\!\varinjlim}_i\operatorname{Cov}(S_i)
\to\operatorname{Cov}(Y)$ is an equivalence if and only if
$\phi_i$ is an isomorphism for some stage. In that event all its
later maps are isomorphisms as well.
\end{statement}
\begin{proof}
The classical covering classifications identify the comparison with
inflation of sets along the maps $\phi_i$. Surjectivity makes every
inflation fully faithful. If the comparison is essentially surjective,
the regular left action of $\pi=\pi_1(Y,y)$ on itself descends to
some stage. Its kernel is trivial, but it must contain
$\ker\phi_i$. Hence $\phi_i$ is injective and therefore an
isomorphism. For a later $j$, the factorization
$\phi_i=(\pi_1(S_j)\to\pi_1(S_i))\phi_j$ and surjectivity of
$\phi_j$ imply that both factors are isomorphisms. Conversely an
eventually constant system of these covering categories has the
stated equivalence. This test involves the regular, usually infinite,
cover and is stronger than agreement on every finite quotient.
\end{proof}

\begin{statement}{F.15}{Proved}{Countable open intersections in a locally compact Hausdorff space}
Let $X$ be locally compact Hausdorff and let $V_n\subset X$ be a
countable family of open subsets. Put $K=\bigcap_nV_n$ with its
subspace topology. The intersection of the associated open sublocales
is spatial and is the locale of $K$. Consequently, for a decreasing
sequence and a closed subset $Z\subset X$,
\[
 \operatorname*{2\!\varprojlim}_n\operatorname{Sh}(V_n\setminus Z)
       \simeq\operatorname{Sh}(K\setminus Z).
\]
This proves Lead 21(A) for manifolds with their usual Hausdorff,
locally compact convention. It does not require $K$ itself to be
closed or compact.
\end{statement}
\begin{proof}
Write $L=\bigcap_n V_n$ as a sublocale of $X$. We prove it has
enough points, including separation of arbitrary opens, rather than
merely computing its point set. Suppose $a,b\in\mathcal O(L)$
satisfy $a\not\leq b$. Choose ambient open representatives $A,B$,
and intersect $L$ with the closed sublocale of $F=X\setminus B$.
In this closed quotient the image of $a$ is still nonzero: the
closed-quotient frame map sends $u$ to $u\vee b$, so its image is
zero precisely when $u\leq b$. The space $F$ is locally compact
Hausdorff.

Let $q:\mathcal O(F)\to\mathcal O(L\cap F)$ be the quotient
frame map. It sends each $V_n\cap F$ to $1$ and sends $A\cap F$
to a nonzero element. Cover $A\cap F$ by relatively compact opens
whose closures lie in $A\cap F$. Since $q$ preserves joins, one
such open $W_0$ has nonzero image. Recursively, after choosing $W_n$,
the open $W_n\cap V_n\cap F$ has the same nonzero image as $W_n$.
Cover it by relatively compact opens with closures inside it, and
choose one $W_{n+1}$ with nonzero image. The nonempty compact sets
$\overline W_n$ are nested in $\overline W_0$, so they have a
common point $x$. By construction $x\in A\cap F$ and
$x\in V_n$ for every $n$. Thus $x\in K\cap A\setminus B$.
As a point of $X$ it factors through every $V_n$, hence through
their localic meet. It separates $a$ from $b$.

Therefore $L$ is spatial. Its points are exactly $K$, and because
the frame map from $X$ is surjective, its point topology is the
subspace topology. Finally apply this result in the locally compact
Hausdorff open subspace $X\setminus Z$, and use F.7's localic
2-limit description. The proof uses countability for the nested
compact construction and does not assert the same spatiality for an
arbitrary intersection of opens.
\end{proof}

\begin{statement}{F.16}{Proved}{The intrinsic stack versions and the fibre convention}
Let $\mathcal X$ be a topological \'etale stack admitting a locally
compact Hausdorff atlas. Let $\mathcal V_n$ be intrinsic open
substacks, decreasing in a countable sequence, and let
$\mathcal Z$ be an intrinsic closed substack. Write
$\mathcal K$ for the restriction of an atlas groupoid to the
intersection of the inverse images of the $\mathcal V_n$.
Then
\[
 \operatorname*{2\!\varprojlim}_n
       \operatorname{Sh}(\mathcal V_n\setminus\mathcal Z)
       \simeq\operatorname{Sh}(\mathcal K\setminus\mathcal Z).
\]
For nested closed substacks $\mathcal N_n$ of a fixed stack
$\mathcal T$, the same statement holds with their closed
intersection, without countability. Applied to intrinsic closed
neighbourhoods in an oriented blow-up whose downstairs intersection
is exactly the residual gerbe of the specified point, the limit is
the inverse image of that gerbe. It need not be the geometric
2-fibre over a chosen object of the stack. The stated shrinking
condition must be imposed: a locally compact Hausdorff atlas alone
does not make the stack's orbit space Hausdorff.
\end{statement}
\begin{proof}
Pull back the diagram along an \'etale surjective atlas. This is
base change of geometric morphisms; the relevant 2-limits commute
with that pullback. On the atlas the first comparison is exactly
F.15. It is compatible with the two pullbacks to the atlas relation
space, since all the opens are intrinsic. Sheaf descent for the
atlas therefore gives the equivalence on the stack. Equivalently,
one computes the open sublocales relative to its topos and checks
their spatial atlas presentations by F.15.

For the closed version, a geometric morphism factors through every
closed subtopos exactly when it kills every complementary open
subterminal. It then kills their union, the complement of the
closed intersection. This is the bicategorical universal property
used in F.9. Under the stated downstairs shrinking condition, the
intersection is the claimed inverse image, and intrinsic
neighbourhoods retain its full isotropy. In $BC_2\times\Delta$
their intersection over the origin
is $BC_2$, whereas base change along a chosen point $*\to BC_2$
removes that isotropy. F.6 computes the resulting difference after
blowing up. An arbitrary shrinking family in a single chart need
not descend to intrinsic open substacks, so it is not covered by
this assertion.
\end{proof}


\begin{statement}{F.17}{Proved}{An exact criterion for every Hausdorff point germ}
Let $S$ be Hausdorff and $x\in S$, with the filter of all open
neighbourhoods. For an open $U\ni x$, define its quasi-component
\[
 Q_U(x)=\bigcap\{A\subset U:A\text{ is clopen in }U, x\in A\}.
\]
The literal inverse sheaf-topos limit is $\operatorname{Set}$.
The fibre comparison
\[
 \operatorname*{2\!\varinjlim}_{U\ni x}\operatorname{Cov}(U)
       \longrightarrow\operatorname{Set}
\]
is an equivalence if and only if $Q_U(x)$ is a neighbourhood of
$x$ in $S$ for every open $U\ni x$. In particular every locally
connected point satisfies the condition. It fails for
$S=\{0\}\cup\{1/n:n\geq1\}\subset\mathbf R$ at $x=0$.
Moreover this condition is equivalent to the existence of any
Grothendieck topos whose locally constant-object category is the
covering-germ category.
\end{statement}
\begin{proof}
For any $z\ne x$, choose disjoint open neighbourhoods $W_z$ of
$z$ and $U_z$ of $x$. A cone over all $\operatorname{Sh}(U)$
pulls $W_z$ to zero, since it factors through $U_z$. The $W_z$
cover $S\setminus\{x\}$, so the cone factors through the
closed point subtopos. This proves the literal limit assertion
using only Hausdorff separation.

Every covering is trivial on some neighbourhood of $x$, so every
object of the covering germ is isomorphic to a constant cover.
Every map between its fibres extends as a constant map; the
comparison is therefore essentially surjective and full.
A map between constant covers with fibres $A,B$ on $U$ is a
family $(f_a:U\to B)_{a\in A}$ of locally constant functions.
The fibre $f_a^{-1}(f_a(x))$ is clopen and contains $x$.
Hence all $f_a$ are simultaneously constant on $Q_U(x)$. If
$Q_U(x)$ contains an open neighbourhood, the germ map equals
its constant fibre map, proving faithfulness and the sufficiency.

Conversely index a constant covering by the set of all clopen
subsets $C$ of $U$ containing $x$, and map its $C$-sheet to the
two-point constant cover by the characteristic function $1_C$.
This map and the constant-$1$ map agree on the fibre at $x$.
Faithfulness requires them to agree on one smaller neighbourhood
$V$; then $V\subset\bigcap_C C=Q_U(x)$. This proves necessity.

For the convergent sequence the singleton $\{1/n\}$ is clopen
for every $n$. Its characteristic function on the $n$th sheet
gives a map from the trivial countable cover to the two-point
cover, zero on the fibre at $0$. Every neighbourhood of $0$
contains some $1/n$, so the map is not the zero germ. Equivalently
$Q_U(0)=\{0\}$ for each sufficiently small $U$, and this is
not a neighbourhood. For a locally connected point, its connected
component in $U$ contains a neighbourhood and lies in $Q_U(x)$,
which verifies the positive assertion.

For the last assertion, if the criterion fails, the characteristic
function family constructed above supplies two different maps from
the constant $I$-cover to $2$ which agree on every standard global
section $1\to I$, where $I$ indexes the clopen sets. Thus that
constant cover is not a coproduct of $I$ copies of $1$. No other
object can be such a coproduct either. Indeed every object of the
point germ is a constant cover up to isomorphism, and its global
sections are exactly its fibre elements: a locally constant section
with a prescribed value is constant on some neighbourhood of $x$.
If a constant $A$-cover with injections indexed by $I$ were a
coproduct, their fibre values would be distinct by testing arbitrary
maps to $2$, and exhaust $A$ by uniqueness for constant maps to $2$.
Hence it would be isomorphic, with the injections, to the same
constant $I$-cover just shown not to be a coproduct. Every category
of locally constant objects in a Grothendieck topos contains all
coproducts of its terminal object. This rules out every such
realization when the criterion fails. When the criterion holds,
the category is $\operatorname{Set}$ itself and has that
realization.
\end{proof}

\begin{statement}{F.18}{Proved}{An exact criterion for residual closed-subset germs}
Let $Y$ be a closed subset of a regular Hausdorff space $X$, and
let $V$ run through a cofinal filter of open neighbourhoods of
$Y$. The comparison
\[
 \operatorname*{2\!\varinjlim}_V\operatorname{Cov}(V\setminus Y)
 \longrightarrow
 \operatorname{LC}\!\left(
     \operatorname*{2\!\varprojlim}_V\operatorname{Sh}(V\setminus Y)
                 \right)
\]
is an equivalence if and only if $Y$ is open in $X$.
\end{statement}
\begin{proof}
F.7 shows that the limit topos is degenerate, whose locally
constant-object category has one object and one arrow. If no
residual stage is empty, its empty and one-sheeted covers remain
nonisomorphic at every stage and therefore in the filtered
category. The comparison is then not an equivalence. An empty
residual stage is an open $V$ with $Y\subset V\subset Y$,
so it exists precisely when $Y$ is open; cofinality includes such
a stage whenever $Y$ is open. At that empty stage every covering
and morphism restricts to the unique one. This proves both
directions, including the empty-subset case.
\end{proof}

\begin{statement}{F.19}{Proved}{Finite covering continuity for every compact Hausdorff germ}
Let $X$ be locally compact Hausdorff, and let $Y\subset X$ be
compact. Over the filter of all open neighbourhoods $V$ of $Y$,
restriction induces an equivalence
\[
 \operatorname*{2\!\varinjlim}_{V\supset Y}
      \operatorname{Cov}_{\mathrm{fin}}(V)
        \simeq \operatorname{Cov}_{\mathrm{fin}}(Y).
\]
Equivalently, finite locally constant objects of the literal limit
$\operatorname{Sh}(Y)$ descend, together with all their morphisms,
to this neighbourhood germ. Neither $Y$ nor $X$ is assumed locally
connected or semilocally simply connected. This holds in particular
for the Cantor example in F.11, where the analogous assertion for
arbitrary coverings is false.
\end{statement}
\begin{proof}
First extend an object. The rank of a finite covering on $Y$ is
locally constant and has finite image by compactness. Split $Y$
into its finitely many compact clopen rank strata and choose
pairwise disjoint ambient neighbourhoods of them. It suffices to
handle a constant rank $n$ cover. Its frame bundle is a principal
bundle for the finite discrete group $G=S_n$.

Choose a finite trivializing open cover of the compact Hausdorff
space $Y$ and a subordinate finite partition of unity. If it has
$m$ members, the partition and the local frames give the usual
explicit equivariant map from the frame bundle to the finite join
\[
                         E_m=G*\cdots*G
\]
of $m$ copies of $G$: the join weights are the partition values
and each nonzero-weight vertex is its local frame coordinate.
The diagonal $G$ action on $E_m$ is free, and the quotient
$B_m=E_m/G$ is a finite polyhedron. The frame bundle is the
pullback of $E_m\to B_m$ by the resulting map $f:Y\to B_m$.
Indeed this comparison is equivariant and is a bijection on
every principal $G$ fibre, hence an isomorphism of coverings.

For clarity, only elementary neighbourhood extension of maps to
a finite polyhedron is needed here. Embed such a polyhedron as
the realization of a finite simplicial complex in a standard
simplex. It is a neighbourhood retract: for simplex coordinates
$p_i$, let $t(p)$ be the maximum, over the minimal nonfaces
$\tau$, of $\min_{i\in\tau}p_i$. On the open neighbourhood
$\max_i p_i>t(p)$, replace $p_i$ by
\[
             \frac{\max(0,p_i-t(p))}
                  {\sum_j\max(0,p_j-t(p))}.
\]
The positive support contains no nonface, and on the complex
$t=0$, so this is a retraction. Positive-part normalization
extends the construction to an open Euclidean neighbourhood.
The case of no nonfaces is the whole simplex, with $t=0$.
A free finite action on a finite polyhedron has a finite
polyhedral quotient: take an invariant sufficiently fine
subdivision so that translated closed stars are disjoint.
This applies to $B_m$ above.

Choose a relatively compact open neighbourhood $U$ of $Y$ in
$X$. Its closure is compact Hausdorff and therefore normal.
Extend the finitely many Euclidean coordinates of $f$ from the
closed subset $Y$ to $\overline U$ by the Tietze extension
theorem. Restrict to the open neighbourhood where the extended
map lies in the retracting neighbourhood, and compose with the
retraction. Pulling back $E_m\to B_m$ and its associated
$n$-element covering extends the original cover. Combining the
finitely many rank strata proves essential surjectivity.

We give the morphism argument separately, since this is exactly
where infinite fibres can fail. A finite covering $p:T\to V$
is proper: the inverse image of a compact set is compact, as
one sees by refining a finite evenly covered open cover of
the compact set to finitely many compact closed pieces. In
particular it is closed. A section of $T|_Y$ has compact image
$S$, and its complement inside $T|_Y$ is also compact. Choose
disjoint open neighbourhoods $O_1,O_2$ of these compact sets in
the Hausdorff space $T$. Since $p$ is closed, remove from $V$
the image of $T\setminus(O_1\cup O_2)$. Over the remaining
neighbourhood of $Y$, the two opens partition the cover.
The number of points in $O_1$ over a base point is locally
constant and is one along $Y$. Restrict again to where it is
one; this degree-one subcover extends the given section.

For two finite coverings $A,B$ on $V$, the fibrewise set of
maps $A_v\to B_v$ is itself a finite covering of $V$, with
its evident local product charts. A morphism over $Y$ is a
section of that covering and therefore extends by the preceding
argument. This proves fullness. If two morphisms agree over
$Y$, their equalizer is a clopen subcover of $A$ containing
$A|_Y$. The projection of its closed complement is closed,
again by properness. Deleting that projection makes the two
morphisms equal, proving faithfulness. Finally F.7 identifies
the literal inverse topos with $\operatorname{Sh}(Y)$.
\end{proof}

